% Asset ID: AS1AlgebraicExpressionsTIKZ-004
% Source: Dr Frost factorising slides and transcript
% Related lesson section: Factorising, Worked Examples, Guided Practice
% Purpose: Show the split-middle-term method for factorising 6x^2+x-2.

\documentclass[tikz,border=8pt]{standalone}
\usetikzlibrary{arrows.meta,positioning}

\begin{document}
\begin{tikzpicture}[
    font=\sffamily,
    arrow/.style={-{Latex[length=3mm]}, thick},
    step/.style={draw, rounded corners=4pt, thick, fill=gray!8, align=left, text width=6.2cm, minimum height=1.15cm},
    final/.style={draw, rounded corners=4pt, thick, fill=black, text=white, align=center, text width=6.2cm, minimum height=1.15cm}
]

\node[font=\sffamily\bfseries\Large, anchor=west] at (-3.3,3.4)
{Factorising by splitting the middle term};

\node[step] (s1) at (0,2.2)
{\textbf{Start}\\$6x^2+x-2$};

\node[step] (s2) at (0,0.9)
{\textbf{Find two numbers}\\Add to $1$ and multiply to $6(-2)=-12$.\\The numbers are $4$ and $-3$.};

\node[step] (s3) at (0,-0.6)
{\textbf{Split the middle term}\\$6x^2+x-2=6x^2+4x-3x-2$};

\node[step] (s4) at (0,-2.1)
{\textbf{Factorise by grouping}\\$6x^2+4x-3x-2$\\$=2x(3x+2)-1(3x+2)$};

\node[final] (s5) at (0,-3.65)
{\textbf{Factor out the repeated bracket}\\$(3x+2)(2x-1)$};

\draw[arrow] (s1.south) -- (s2.north);
\draw[arrow] (s2.south) -- (s3.north);
\draw[arrow] (s3.south) -- (s4.north);
\draw[arrow] (s4.south) -- (s5.north);

\node[draw, rounded corners=4pt, thick, fill=gray!12, align=left, text width=6.2cm] at (0,-5.1)
{\textbf{Check:} $(3x+2)(2x-1)=6x^2-3x+4x-2=6x^2+x-2$.};

\end{tikzpicture}
\end{document}
